\section{Entropy increments in moving frames}
\label{sec:entropy}

The stationary packet masses determine how many event--matrix cells are
occupied at each scale.  We express the change in this count through the
information gained or lost in the four matrix coordinates.  Their widths
change at different rates, so each coordinate must be measured conditional
on the coordinates already observed.  The resulting entropy demand will
be compared with the planar projection constraints in
Section~\ref{sec:projections}.

The use of entropy increments to record projection growth across scales
is related to the local entropy averages method of \citet{HS2012}.
Here we use finite conditional entropies in the stationary moving frames.

\subsection{The limiting entropy function}

Use the weighted event law of Proposition~\ref{prop:stationary-data}.
Each retained collapsed pair $(i,I_1)$ has weight $\omega_i$, and
probabilities are obtained by dividing by the total weight of these pairs.
Thus the index marginal is proportional to $\omega_i$ times the number of
its retained pairs; conditional on the index, those pairs are equally
likely.  The matrix belonging to the sampled index is denoted by $M$.
All finite entropies in this section are Shannon entropies, with natural logarithms.
We write
\[
 \mathrm H_N(Z\mid W)=\frac{\mathrm H(Z\mid W)}{\log N},
 \qquad
 \mathrm I_N(Z;Z'\mid W)=\frac{\mathrm I(Z;Z'\mid W)}{\log N}.
\]
After taking the subsequences below, $\mathrm H_\infty$ and
$\mathrm I_\infty$ denote their limits.  Changing any observation by a
partition of bounded mutual overlap changes its ordinary entropy by
$O(1)$, and hence changes its limiting normalized entropy by zero.

Scalar observation at threshold $r$ means quantization in a nested
dyadic grid at length $N^{-r}$, with the integer dyadic depth rounded.
Recall that $X,U$ are the normal-position and normal-velocity forms,
while $Y,V$ are the second position and velocity components in the fixed
spatial chart.  In particular, $V$ is an absolute velocity observation,
unchanged by the event frame.  The packet widths
\eqref{eq:estimate-widths} give the thresholds
\begin{equation}
 (k_X,k_Y,k_U,k_V)=(0,\tau,\ell,1),\qquad
 r_i(c,a)=\rho+(w-k_i)c+a.
 \label{eq:entropy-thresholds}
\end{equation}
Increasing the offset $a$ refines all four coordinates equally. Increasing
$c$ changes the threshold for coordinate $i$ at rate $w-k_i$.
Let $Q(c,a)$ be their joint observation at these thresholds, in the frame
of a fixed representative of the event label $e_c$.  Representatives can
be bin centers.  On every compact interval of $c\in(0,1)$, all four
thresholds lie strictly between $0$ and $1$ for sufficiently small
positive offsets $a$.  Negative offsets are also allowed when $w<1$.
More explicitly, the domain of interior full-cell tests is
\begin{equation}
 \mathcal U=\{(c,a):0<c<1,
       -\rho(1-c)<a<w(1-c)\}.
 \label{eq:entropy-domain}
\end{equation}
When $w=1$, offset zero is the left boundary of this domain.

Apply Lemma~\ref{lem:regularization} also to joint observations from a
countable dense family of $c$'s, offsets, extra label depths, and full or
partial coordinate observations.  Include unequal-depth observations
and the preceding full cells.  Their occupied joint cells have common
mass exponents, and their conditional refinement counts have common
exponents.  Consequently a limiting conditional entropy is also the
limiting logarithm of the corresponding typical conditional count.
This is why later entropy identities yield covering and ball-mass
estimates.  Mean entropy alone would not give those conclusions.
At every finite stage there are only finitely many tests.  A new finite
list can be added before the next sufficiently large value of $N$ is
chosen.  The regularization and deficiency deletions preserve all earlier
limits and the stationary packet data.  We use this slow diagonal
throughout the section.

Here are the comparison estimates that extend the limiting full-cell
tests to continuous parameters.  Changing the frame by
$(\Delta t,\Delta\vartheta)$ changes the forms, including their values
on matrix differences, according to
\begin{align}
 V'&=V,& U'&=U+\Delta\vartheta V,&
 Y'&=Y+\Delta t V,\nonumber\\
 X'&=X+\Delta t U+\Delta\vartheta Y
                      +\Delta t\Delta\vartheta V.
 \label{eq:entropy-frame-change}
\end{align}
If both frame labels refine the same coarse label $e_c$, then
$|\Delta t|\lesssim N^{-\ell c}$ and
$|\Delta\vartheta|\lesssim N^{-\tau c}$.  These factors are exactly
the ratios of the widths in \eqref{eq:entropy-thresholds}.  Conditional
on a common finer label, each full cell in either frame therefore meets
boundedly many full cells in the other frame.  Changing thresholds or
label precisions by at most $\epsilon$ costs
$O(\epsilon)+o(1)$ in normalized entropy: the extra label and grid
refinements have at most $N^{O(\epsilon)}$ possibilities.

A partial coordinate test requires a slightly stronger comparison.
Start inside a full cell and give both frames extra label precision
$\sigma>0$.  For a refinement gap of size $b$ with
$b\ll\min\{\ell,\tau\}\sigma$, errors from the unobserved
coordinates in \eqref{eq:entropy-frame-change} are smaller than the
requested widths.  Known coordinates cause only known translations,
which do not affect intersection counts.  Partial rectangular tests
then have the same stable comparison under nearby parameters.  All
continuous partial tests below are used in this regime.  The limit in
$N$ is taken first, with $\sigma$ fixed, then the derivative gap tends
to zero, and only afterward does $\sigma$ tend to zero.

The stable full-cell comparisons define a locally Lipschitz function
\begin{equation}
 \mathscr H(c,a)=
       \lim_{N\to\infty}\mathrm H_N(e_c,Q(c,a)).
 \label{eq:entropy-function}
\end{equation}
At offset zero, each saturated packet occupies only subpower many joint
event--matrix cells.  The local upper mass bound gives the reverse count,
so the stationary mass formula determines their entropy:
\begin{equation}
 \mathscr H(c,0)=\lim_{N\to\infty}\log_N(n/n_c)+s_0\ell c,
 \qquad
 \partial_c\mathscr H(c,0)=dw-B_*-\tau+s_0\ell.
 \label{eq:entropy-central-value}
\end{equation}
To see both directions of the first equality, a full matrix cell has
active index mass at most $n_cN^{o(1)}$ by the local bound, and a time
bin at scale $c$ has relative mass at most
$N^{-s_0\ell c+o(1)}$ on each index.  Thus each occupied joint cell
has event probability at most
$(n_c/n)N^{-s_0\ell c+o(1)}$.  Conversely, the disjoint assigned
packets at each time block have mass at least $n_c$ in exponent.
Summing their assigned masses, and then summing the regular time counts,
gives at most $(n/n_c)N^{s_0\ell c+o(1)}$ occupied joint cells.
Each packet splits among only subpower many corresponding grid cells
and normal labels.  The two estimates prove the equality.  Differentiating
\eqref{eq:stationary-packet-mass} proves its derivative formula.

The same local bound applied to cells whose four widths are multiplied
by $N^{-a}$ gives
\begin{equation}
 \mathscr H(c,a)\ge\mathscr H(c,0)+da
 \label{eq:entropy-offset-bound}
\end{equation}
whenever the tested widths remain in range.  This uses their total
homogeneity $x+y=d$.  It applies on both sides of zero if $w<1$,
and on the positive side if $w=1$.  Write
$g=\partial_a\mathscr H$, defined almost everywhere.

\subsection{A differentiation lemma for entropy gaps}

We next express the change in $\mathscr H$ as $c$ increases.  A scale
increment $b$ adds at least $s_0\ell b$ of time information and changes
the four matrix thresholds by $(w-k_i)b$.  We will construct conditional
rates $j_G$ for groups $G$ of equal-speed coordinates and prove
\[
 \partial_c\mathscr H\ge s_0\ell+\sum_G(w-k_G)j_G.
\]
Here $j_G$ measures the additional information in $G$ after the groups
with smaller $k_i$ have been observed.  Equal speeds must be kept together
because their refinements occur at the same depths.

Two issues must be resolved before this calculation is valid.  The
coordinate thresholds move by different amounts, so we need a first-order
formula for rectangles with unequal depths, uniformly over all intervals
in a shrinking neighborhood.  Also, a change of frame mixes coordinates;
we first condition on a finer event label to control this mixing and then
remove that extra information.  The next lemma supplies the uniform
differentiation needed for the rectangular observations.

\begin{lemma}[Uniform differentiation of a superadditive gap kernel]
\label{lem:gap-density}
Let $K(s,t)$, $s<t$ in an interval, be a nonnegative kernel satisfying
\[
 K(s,u)\ge K(s,t)+K(t,u),\qquad
 0\le K(s,t)\le C(t-s).
\]
There is a measurable $\phi$, $0\le\phi\le C$, such that for almost
every $a$, and for every fixed $C_0<\infty$,
\begin{equation}
 \sup_{[s,t]\subset[a-C_0b,a+C_0b]}
 \bigl|K(s,t)-\phi(a)(t-s)\bigr|=o(b)
       \quad(b\downarrow0).
 \label{eq:gap-uniform-density}
\end{equation}
\end{lemma}

\begin{proof}
For a finite partition $\mathcal P$ of $[s,t]$, sum the values of $K$
on its intervals, and set
\[
 m(s,t)=\inf_{\mathcal P}\sum_{[u,v]\in\mathcal P}K(u,v).
\]
Concatenation of partitions gives
$m(s,u)\le m(s,t)+m(t,u)$.  Inserting $t$ into any partition of
$[s,u]$ can only decrease its sum, by superadditivity of $K$.
Taking infima gives the reverse inequality.  Thus $m$ is additive and
$0\le m(s,t)\le C(t-s)$.  It is the interval function of an
absolutely continuous measure, so
$m(s,t)=\int_s^t\phi(r)\,dr$ for some $0\le\phi\le C$.

The excess $E=K-m$ is nonnegative and superadditive.  Its infimum of
partition sums is zero on every interval.  Fix a compact interval $J$
and $\epsilon>0$.  Consider points in its interior that admit arbitrarily
small centered intervals $I$ with $E(I)>\epsilon|I|$.
Choose a finite partition of $J$ with sum of excesses less than an
arbitrary $\eta>0$.  Except at its finitely many endpoints, the small
test intervals can be confined to individual partition cells.  The
Vitali covering theorem supplies a disjoint countable selection of
these intervals covering the exceptional set up to a null set.
For every finite selection, superadditivity and nonnegativity give
\[
 \epsilon\sum_I|I|<\sum_I E(I)
             \le\sum_{J'\in\mathcal P}E(J')<\eta.
\]
Pass to the countable union.  Its length, and hence the measure of the
exceptional set, is at most $\eta/\epsilon$.  Since $\eta$ was
arbitrary, this set is null.  Exhaust the domain by compact intervals
and take countably many positive rational $\epsilon$.  It follows that
the excess on centered intervals has upper density zero almost
everywhere.

At a point with this property that is also a Lebesgue point of $\phi$,
the excess on any subinterval of $[a-C_0b,a+C_0b]$ is bounded by the
excess on the entire surrounding interval.  The additive part has error
at most
\[
 \int_{a-C_0b}^{a+C_0b}|\phi(r)-\phi(a)|\,dr=o(b).
\]
This proves \eqref{eq:gap-uniform-density}, first for positive integral
$C_0$ and therefore for every fixed $C_0$.
\end{proof}

Fix $c$ and an extra label depth $\sigma>0$ with $c+\sigma<1$.
Work in the $e_{c+\sigma}$ frame, and condition on that label.
For an axis subset $S$, let $S_t$ denote its observation at offset $t$
and let $Q_s$ denote the full observation at offset $s$, both with the
threshold base $c$.  Define
\begin{equation}
 K_S^\sigma(s,t)=\mathrm H_\infty
          (S_t\mid Q_s,e_{c+\sigma}),\qquad s<t.
 \label{eq:entropy-gap-kernel}
\end{equation}
Use only a short offset interval compared with $\sigma$, as above.
The kernel is bounded between zero and $|S|(t-s)$, and is
superadditive.  Indeed, split the entropy at $S_t$ in the entropy
of $S_u$ given $Q_s$; then condition the second increment further
on the remaining coordinates of $Q_t$.  Lemma~\ref{lem:gap-density}
gives densities $\phi_S^\sigma(c,a)$ with the uniform expansion
\eqref{eq:gap-uniform-density}.

This yields a rectangular layering rule.  Start with a full observation
at $a_0$.  To reveal a rectangular refinement, move through the finitely
many depths at which its set of required axes changes.  On a layer
$[d_1,d_2]$, call the required axes $S$.  The existing information
contains the full starting cell and $S_{d_1}$, and is contained in
the full observation at $d_1$.  Therefore the entropy contributed on
that layer is between
\begin{equation}
 K_S^\sigma(d_1,d_2)\quad\hbox{and}\quad
 K_S^\sigma(a_0,d_2)-K_S^\sigma(a_0,d_1).
 \label{eq:rectangular-layer-bounds}
\end{equation}
If all depths are within $O(b)$ of a differentiation point $a$, both
bounds equal $\phi_S^\sigma(c,a)(d_2-d_1)+o(b)$.
One may sum the finitely many layers, or subtract two such rectangular
entropy formulas.  The latter operation permits nuisance coordinates to
be conditioned upon at a different depth.  The error is uniform over
the endpoints of these rectangles in a fixed $O(b)$ window.

\subsection{Removing the extra event precision}

The total information in an extra label of depth $\sigma$ is
$O(\sigma)$.  This does not control an entropy increment divided by its
length $b$ when $b$ tends to zero first.  We instead bound the extra
label's information on the same shrinking interval.  Its density will
tend to zero as $\sigma$ decreases.

\begin{lemma}[Local removal of label conditioning]
\label{lem:remove-label}
Let $\sigma\downarrow0$ through a fixed countable sequence of extra
label depths.  For almost every $(c,a)$, all the rates
$\phi_S^\sigma(c,a)$ converge to rates $\phi_S(c,a)$.
The rectangular layering rule holds with these limiting rates, with
the order of limits stated above.  Moreover
\begin{equation}
 \phi_\varnothing=0,\qquad
 \phi_{\{X,Y,U,V\}}=g.
 \label{eq:entropy-full-rate}
\end{equation}
For two nested extra labels, changes in conditional increments are
controlled by the incremental mutual information of a full observation
over a containing gap.  Its density tends to zero almost everywhere as
the total extra label depth tends to zero.
\end{lemma}

\begin{proof}
Use first the fixed $e_c$ frame and define
\[
 F_\sigma(a)=\mathrm I_\infty
                 (Q(c,a);e_{c+\sigma}\mid e_c).
\]
Nested grids and the chain rule show that $F_\sigma$ is nondecreasing,
and that its increment on $[a_1,a_2]$ is exactly
\begin{equation}
 F_\sigma(a_2)-F_\sigma(a_1)
  =\mathrm I_\infty
    (Q(c,a_2);e_{c+\sigma}\mid Q(c,a_1),e_c).
 \label{eq:label-increment-information}
\end{equation}
It is bounded by $4(a_2-a_1)$, and its total variation on any fixed
interior interval is at most
$\mathrm H_\infty(e_{c+\sigma}\mid e_c)=O(\sigma)$.
If $\sigma_1\le\sigma_2$, the increment measure for $\sigma_1$
is at most that for $\sigma_2$, again by the chain rule and label
nesting.  Their bounded densities therefore have decreasing versions
along $\sigma\downarrow0$.  Their integrals tend to zero, so their
density limit is zero almost everywhere.  This monotonicity is the
reason that the conclusion is almost everywhere, rather than just in
$L^1$.

We spell out how to use this measure for partial observations, since
mutual information need not decrease under arbitrary conditioning.
Write $L_0=e_c$ and take nested extra labels
$L_1=e_{c+\sigma_1}$ and $L_2=e_{c+\sigma_2}$, where
$0<\sigma_1\le\sigma_2$. First use the $L_1$ frame for every
partial observation. Let $A$ consist of the full starting cell and any
preceding nuisance observations, and let $Z$ be the partial increment
being revealed. Choose containing full $L_0$-frame observations
$R=Q(c,a-Cb)$ and $T=Q(c,a+Cb)$.
Given $L_1$, the observation $A$ determines $R$, and $(T,L_1)$
determines $(A,Z)$, up to bounded overlap. Adjoin the bounded-overlap
observations at ordinary entropy cost $O(1)$. The chain rule gives
\begin{align}
 \mathrm I(Z;L_2\mid A,L_1)
 &\le \mathrm I(T;L_2\mid R,L_1)+O(1)\nonumber\\
 &\le \mathrm I(T;L_2\mid R,L_0)+O(1).
 \label{eq:label-containing-chain}
\end{align}
For the second inequality use $L_1\subset L_2$: its difference is
the nonnegative term $\mathrm I(T;L_1\mid R,L_0)$.
After dividing by $\log N$ and taking the limit, the last line is
$F_{\sigma_2}(a+Cb)-F_{\sigma_2}(a-Cb)$.
The first line bounds exactly the change in
$\mathrm H(Z\mid A,L_1)$ on adding $L_2$. Thus nuisance data are
permitted because they are determined by the containing full ending
grid and the older label; no arbitrary-conditioning assertion is used.
A difference of rectangular entropies is treated by applying this
bound to its finitely many terms.

Only after conditioning on $L_2$ do we change the partial observations
to the $L_2$ frame. Full starting cells in the two frames have bounded
mutual overlap. For partial refinements retain this full base cell and
take $b\ll\min\{\ell,\tau\}\sigma_1$. The unresolved base
coordinates in \eqref{eq:entropy-frame-change} acquire the extra
factors $N^{-\ell\sigma_1}$ or $N^{-\tau\sigma_1}$, which put
all their errors inside the ending widths. Known coordinate centers
cause only known translations. The joint partial tests therefore have
bounded mutual overlap given $L_2$ and the base cell. This comparison
costs $O(1)$ ordinary entropy, which vanishes in the $N$ limit before
any division by $b$.

Apply Lemma~\ref{lem:gap-density} and differentiate at a common
Lebesgue point of the information densities.  The difference of two
$\sigma$-rates is bounded by a constant times the information density
for their total label refinement.  That density tends to zero almost
everywhere.  The rates are thus Cauchy there, proving convergence and
the limiting layering rule.  For the full subset, conditional
increments differ from increments of $\mathscr H(c,a)$ by exactly
the information in \eqref{eq:label-increment-information}.  This proves
\eqref{eq:entropy-full-rate}.  Taking a countable intersection for
subsets, extra depths, and stable tests, and then using Fubini in $c$,
gives the assertions almost everywhere in $(c,a)$.  The stable
comparison estimates give measurable versions of all the functions
used here.
\end{proof}

Order the coordinates by increasing speed $k_i$, keeping tied speeds
together.  Let $G_1,\ldots,G_m$ be these groups and
$S_j=G_1\cup\cdots\cup G_j$, with $S_0=\varnothing$.  Define
\begin{equation}
 j_{G_j}=\phi_{S_j}-\phi_{S_{j-1}},\qquad
 \mathcal M=\sum_{j=1}^m k_{G_j}j_{G_j}.
 \label{eq:entropy-group-rates}
\end{equation}
The conditioning in these rates can be read directly from the chain
rule.  In the $e_{c+\sigma}$ frame, for $r>0$ sufficiently small,
\begin{align}
 &K_{S_j}^\sigma(a,a+r)-K_{S_{j-1}}^\sigma(a,a+r)\nonumber\\
 &\qquad=\mathrm H_\infty\bigl((G_j)_{a+r}\mid
             Q_a,(S_{j-1})_{a+r},e_{c+\sigma}\bigr).
 \label{eq:entropy-group-conditioning}
\end{align}
Dividing by $r$, then taking $r\downarrow0$ and finally
$\sigma\downarrow0$, gives $j_{G_j}$.  Thus this rate measures new
information in $G_j$ after the full old cell and the refinements of all
earlier groups are known.  In particular,
\begin{equation}
 0\le j_{G_j}\le|G_j|,\qquad \sum_j j_{G_j}=g.
 \label{eq:entropy-rate-bounds}
\end{equation}
In every case write $a_1=j_X$ and $v_1=j_V$.
For distinct speeds put $y_1=j_Y$ and $u_1=j_U$.
If $\tau=\ell=1/2$, keep the middle pair together and denote its
rate by $m_1=j_{\{Y,U\}}$. In every case the last group is the single axis $V$.

The geometric estimate from Section~\ref{sec:stationary-estimates} now
has a direct interpretation in terms of the first rate.

\begin{corollary}[Normal-position rate]\label{cor:normal-position-rate}
For almost every $(c,a)$ with $0<c<1$ and $0<a<w(1-c)$,
\begin{equation}
 a_1(c,a)\le x=d-1-h<1.
 \label{eq:entropy-normal-position-rate}
\end{equation}
\end{corollary}

\begin{proof}
Fix an extra label depth $\sigma>0$ with $c+\sigma<1$.
Lemma~\ref{lem:conditional-X} bounds the normalized conditional
$X$ entropy at every sufficiently small fixed gap $r$ by $xr$ in the
limit $N\to\infty$.  Its base cell is $Q(c,a)$ in the $e_c$ frame.
Given $e_{c+\sigma}$, this full cell and the full cell $Q_a$ in the
finer frame determine one another up to bounded overlap, by
\eqref{eq:entropy-frame-change}.  Replacing the base cell therefore
changes the ordinary conditional entropy by $O(1)$, which vanishes in
the $N$ limit before division by $r$.  The resulting entropy is exactly
$K_{\{X\}}^\sigma(a,a+r)$.  Dividing by $r$ and taking
$r\downarrow0$ gives $\phi_{\{X\}}^\sigma\le x$ at almost every
point.  Finally Lemma~\ref{lem:remove-label}, with
$\sigma\downarrow0$, gives $a_1=\phi_{\{X\}}\le x$.
The positive-offset range is the range required by
Lemma~\ref{lem:conditional-X}.
\end{proof}

\Needspace{10\baselineskip}
\begin{proposition}[Entropy demand]
\label{prop:entropy-demand}
At almost every point of $\mathcal U$,
\begin{equation}
 \partial_c\mathscr H(c,a)\ge s_0\ell+wg-\mathcal M.
 \label{eq:entropy-demand}
\end{equation}
\end{proposition}

\begin{proof}
For $b>0$, label nesting gives the exact normalized limiting identity
\begin{align}
 \mathscr H(c+b,a)-\mathscr H(c,a)
 ={}&\mathrm H_\infty(e_{c+b}\mid e_c,Q(c,a))\nonumber\\
 &+\mathrm H_\infty(Q(c+b,a)\mid e_{c+b})
       -\mathrm H_\infty(Q(c,a)\mid e_{c+b}).
 \label{eq:entropy-demand-chain}
\end{align}
Condition the first entropy additionally on the index.  The old
matrix cell is then deterministic given that index and $e_c$.
The time branching in \eqref{eq:stationary-time-branching} contributes
at least $s_0\ell b$.  The normal range on an index in the old time
bin has only subpower many coarse angular choices, so conditioning on
the old joint label does not cost a positive exponent of that time
branching.  The finite regularization includes these conditional label
tests.  This proves the stated lower bound for the first term.

Fix now $\sigma>b$ and condition both remaining terms further on
$e_{c+\sigma}$.  Given $e_{c+b}$, the two full observations are
sandwiched between full observations at offsets $a-Cb$ and $a+Cb$
in the $c$ frame, up to bounded overlap; all four speeds are bounded.
The change in their entropy difference is bounded by the extra-label
information on this containing gap, as proved in
Lemma~\ref{lem:remove-label}.  At a common differentiation point it
is at most $b$ times a quantity tending to zero as $\sigma\downarrow0$,
plus $o(b)$.

In the common finer frame, changing $c$ to $c+b$ changes the four
thresholds by $(w-k_i)b$.  Use a common full starting offset below
all the thresholds, and subtract the two rectangular layering
formulas.  The axes with smaller $k_i$ extend farther.  The contribution
is therefore
\[
 b\sum_G(w-k_G)j_G^\sigma+o(b).
\]
Let $b\downarrow0$ with $\sigma$ fixed, and then let
$\sigma\downarrow0$.  Combining this expression with the label
contribution gives \eqref{eq:entropy-demand}.
\end{proof}

\subsection{A strong trace for the last coordinate}

The projection constraints will distinguish the sets where $v_1$ is
positive or zero.  Weak convergence alone does not preserve that
distinction as the offset tends to zero.  We therefore prove strong
compactness for $v_1$, using the fact that $V$ is unchanged by frame
changes.

Keep its absolute threshold fixed.  As $c$ increases, the label and
the observations of the other three coordinates become finer.  We will
show that these later data determine the earlier conditioning up to
bounded overlap, while the $V$ increment itself remains the same.  Its
conditional entropy rate is therefore nonincreasing along each fixed
threshold, which gives the required compactness.

\begin{lemma}[Velocity rate along a fixed threshold]
\label{lem:velocity-trace}
The function $v_1(c,a)$ has an essentially nonincreasing version as
$c$ increases on lines where
\[
 \lambda=\rho+(w-1)c+a
\]
is fixed.  If $w<1$, it has strong one-sided traces in
$L^1_{\mathrm{loc}}(dc)$ at $a=0$, through a full-measure set of
offsets.  Its one-sided stripe averages have the same strong limits.
If $w=1$, every sequence of such generic positive offsets tending to
zero has a subsequence on which $v_1(\cdot,a)$ converges strongly in
$L^1_{\mathrm{loc}}(dc)$.
\end{lemma}

\begin{proof}
At fixed $\lambda$, the last conditional rate, before removing an
extra label, is the derivative at $r=0+$ of
\begin{equation}
 \mathrm H_\infty\bigl(V_{\lambda+r}\mid
    V_\lambda,X_{\lambda+c+r},Y_{\lambda+\ell c+r},
    U_{\lambda+\tau c+r},e_{c+\sigma}\bigr).
 \label{eq:velocity-conditional-rate}
\end{equation}
This follows by subtracting the rectangular entropy for the first
three axes from that for all four, using
\eqref{eq:rectangular-layer-bounds}.  Compare two values of $c$ with
the same $\lambda$, using extra precision $\sigma$ at both locations.
The later label refines the earlier one.  The later other-coordinate
thresholds are also finer.  In converting from the later frame to the
earlier frame, the unknown part of $V$ beyond $V_\lambda$ enters
$Y,U,X$ with factors bounded respectively by
\[
 O(N^{-\ell(c+\sigma)}),\qquad
 O(N^{-\tau(c+\sigma)}),\qquad O(N^{-(c+\sigma)}),
\]
where $c$ is the earlier location.  These errors fit in the earlier
end cells when $r\ll\min\{\ell,\tau\}\sigma$.
For clarity, write the two locations as $c_1<c_2$ and put
$\Delta c=c_2-c_1$. In the conversion
\[
 X_1=X_2-\Delta t\,U_2-\Delta\vartheta\,Y_2
                         +\Delta t\Delta\vartheta V,
\]
the four uncertainty terms, divided by the earlier $X$-cell width
$N^{-\lambda-c_1-r}$, are bounded respectively by constants times
\[
 N^{-\Delta c},\quad
 N^{-\tau\Delta c-\ell\sigma},\quad
 N^{-\ell\Delta c-\tau\sigma},\quad N^{r-\sigma}.
\]
The unresolved $V$ terms in $Y_1$ and $U_1$, divided by their
respective earlier widths, are bounded by
$O(N^{r-\ell\sigma})$ and $O(N^{r-\tau\sigma})$.
The remaining observed-coordinate errors are smaller. Hence the later conditioning
data determine the earlier conditioning data up to bounded overlap.
The variable $V$ itself is identical in both frames.  Conditional
entropy decreases when information is added, proving the required
inequality for \eqref{eq:velocity-conditional-rate}.  First differentiate
in $r$ and then remove the extra label by
Lemma~\ref{lem:remove-label}.  Fubini in the common threshold and the
two locations gives the essential monotonicity assertion. More explicitly,
the change of variables $(c,a)\mapsto(c,\lambda)$ is an invertible
shear, and its interior domain is $c>0$, $\lambda>0$, $c+\lambda<1$.
Rate-existence and label-removal null sets remain null in these variables.
After intersecting over countably many extra precisions, Fubini gives,
for almost every $\lambda$, the comparison for almost every ordered
pair $c_1<c_2$. Integrating over two ordered small neighborhoods of
Lebesgue points and shrinking them gives the comparison at every pair
of Lebesgue points on that line. Extend by one-sided limits to obtain
a monotone version. Taking limits of one-sided local averages supplies
a jointly measurable choice of these versions.

For the trace statement suppose $\alpha_0=1-w>0$, and put
$f(c,a)=v_1(c,a)$.  On a compact interior strip define
\[
 h(r,a)=f(r+a/\alpha_0,a).
\]
Here the common threshold is $\rho-\alpha_0r$.  For almost every $r$,
the function $a\mapsto h(r,a)$ has a bounded nonincreasing version,
and hence has one-sided limits $h_+(r)$ and $h_-(r)$ at zero.
Dominated convergence gives strong $L^1$ convergence on compact
$r$ intervals on either side.  For $J\Subset J'\Subset(0,1)$,
undoing the shear yields, for small $|a|$ on the appropriate side,
\begin{align*}
 \|f(\cdot,a)-h_\pm\|_{L^1(J)}
 \le{}&\|h(\cdot,a)-h_\pm\|_{L^1(J')}\\
      &+\|h_\pm(\cdot-a/\alpha_0)-h_\pm\|_{L^1(J)}\longrightarrow0.
\end{align*}
The last term tends to zero by continuity of translation in $L^1$.
Fubini gives one full-measure set of offsets whose original slices
agree almost everywhere with these representatives.  Thus convergence
holds through every sequence in that set.  Averaging the $L^1$
estimate over a one-sided stripe gives the same limit for stripe
averages.

If $w=1$, the fixed-threshold lines are constant-offset lines.
Every generic slice is a bounded monotone function of $c$.  Helly's
selection theorem, followed by dominated convergence, gives strong
$L^1$ compactness on each compact $c$ interval.  A compact exhaustion
and subsequence diagonal give strong $L^1_{\mathrm{loc}}$ compactness.
No choice of values on exceptional offset slices is needed in either
case.
\end{proof}

\subsection{Demand on thin stripes}

\begin{lemma}[Stripe limits of the demand]
\label{lem:stripe-demand}
Fix $J\Subset(0,1)$.  For sufficiently small $\epsilon>0$, the right
stripe averages of $g$ are at least $d$ almost everywhere on $J$:
\begin{equation}
 \frac1\epsilon\int_0^\epsilon g(c,a)\,da\ge d.
 \label{eq:entropy-right-stripe}
\end{equation}
If $w<1$, the left stripe averages satisfy
\begin{equation}
 \frac1\epsilon\int_{-\epsilon}^0g(c,a)\,da\le d.
 \label{eq:entropy-left-stripe}
\end{equation}
Along any sequence of generic admissible offsets tending to zero,
\begin{equation}
 \partial_c\mathscr H(\cdot,a)
    \overset{*}{\rightharpoonup}dw-B_*-\tau+s_0\ell
           \quad\hbox{in }L^\infty(J).
 \label{eq:entropy-boundary-derivative}
\end{equation}
Consequently every weak-star subsequential limit of
\begin{equation}
 \mathcal M-w(g-d)
 \label{eq:entropy-adjusted-demand}
\end{equation}
is at least $\tau+B_*$ almost everywhere on $J$.  This conclusion
holds both for generic offset slices and for one-sided stripe averages.
All these limit inequalities may be tested against arbitrary
nonnegative $L^1(J)$ functions.
\end{lemma}

\begin{proof}
The fundamental theorem of calculus in the offset variable and
\eqref{eq:entropy-offset-bound} give
\[
 \mathscr H(c,\epsilon)-\mathscr H(c,0)
           =\int_0^\epsilon g(c,a)\,da\ge d\epsilon.
\]
On the left, that same inequality at $a=-\epsilon$ gives
$\mathscr H(c,0)-\mathscr H(c,-\epsilon)\le d\epsilon$.
This proves the two stripe bounds; they are not assertions about
individual offset slices of $g$.

Local Lipschitz continuity gives uniform convergence
$\mathscr H(\cdot,a)\to\mathscr H(\cdot,0)$ on $J$ and a common
bound for the $c$ derivatives.  Against a smooth compactly supported
test function $\psi$, integration by parts therefore gives
\[
 \int_J\psi(c)\partial_c\mathscr H(c,a)\,dc
 =-\int_J\psi'(c)\mathscr H(c,a)\,dc
 \longrightarrow
 (dw-B_*-\tau+s_0\ell)\int_J\psi(c)\,dc.
\]
The common $L^\infty$ bound and $L^1$ approximation extend this to
all $L^1$ tests, proving \eqref{eq:entropy-boundary-derivative}.
It also proves the corresponding statement for stripe averages.
Finally Proposition~\ref{prop:entropy-demand} rearranges as
\[
 \mathcal M-w(g-d)\ge s_0\ell+wd-\partial_c\mathscr H.
\]
The right side converges weak-star to $\tau+B_*$.  Nonnegative
$L^1$ tests preserve the inequality under every weak-star
subsequential limit, proving the claim.
\end{proof}

The distinction between weak stripe bounds and the strong trace of
$v_1$ will matter in the closing argument.  Lemma~\ref{lem:velocity-trace}
allows localization to the sets where the limiting velocity rate is
positive or zero, while Lemma~\ref{lem:stripe-demand} supplies the
same lower demand on those measurable sets.
